\documentclass{article}
\usepackage[UTF8]{ctex}
\usepackage{amsmath,amssymb,amsfonts}
\usepackage{enumitem}
\usepackage{geometry}
\geometry{a4paper, margin=1in}
% === 错题本排版优化（不参与题目标识） ===
\setlength{\parindent}{0pt}
\setlength{\parskip}{0.35em plus 0.10em minus 0.08em}
\setlength{\emergencystretch}{3em}
\setlist{itemsep=0.25em, topsep=0.40em, parsep=0pt}
\allowdisplaybreaks[2]
% === 错题本排版优化结束 ===
\begin{document}
\section*{上海财经大学2024年数学分析真题}
\section*{50．上海财经大学2024年数学分析试卷}

\subsection{题目1}
1．计算 \(\lim _{x \rightarrow+\infty}\left(x^{\frac{2}{x}}-1\right)^{\frac{2}{\ln x}}\) ．

\subsection{题目2}
2．计算 \(\lim _{n \rightarrow \infty} \int_{0}^{1} \frac{(2 n-1) x^{n-1}}{1+2^{x}} \mathrm{~d} x\) ．

\subsection{题目3}
3．设 \(f(x)=\left\{\begin{array}{ll}\sin \frac{x}{2}, & x>0 ; \\ x \arctan x, & x \leq 0 .\end{array}\right.\) 求 \(\int_{0}^{\pi+1} f(x-1) \mathrm{d} x\) ．

\subsection{题目4}
4．解答如下问题：

\begin{itemize}
  \item[（1）] 计算 \(f(x)=\frac{1}{1-x-x^{2}}\) 在 \(x=0\) 处的幂级数展开式。
  \item[（2）] 设 \(a_{0}=a_{1}=1, a_{n}=a_{n-2}+a_{n-1}(n \geq 2)\) ，求幂级数 \(\sum_{n=0}^{\infty} a_{n} x^{n}\) 的和函数．
\end{itemize}

\subsection{题目5}
5．求一曲线旋转得到的旋转体的体积。

\subsection{题目6}
6．求曲线积分 \(\int_{L} y^{3} \mathrm{~d} x+\left|3 x^{2}-y\right| \mathrm{d} y\) ，其中 \(L: x^{2}+y^{2}=4\) ，取逆时针方向．

\subsection{题目7}
7．求曲面积分

\[
\iint_{\Sigma}\left(y^{2}-z^{2}\right) \mathrm{d} y \mathrm{~d} z+\left(z^{2}-x^{2}\right) \mathrm{d} z \mathrm{~d} x+\left(x^{2}-y^{2}\right) \mathrm{d} x \mathrm{~d} y
\]

其中 \(\sum\) 为上半椭球面 \(\frac{x^{2}}{a^{2}}+\frac{y^{2}}{b^{2}}+\frac{z^{2}}{c^{2}}=1, z \geq 0(a, b, c>0)\) ，方向取上侧．

\subsection{题目8}
8．已知 \(\left(\frac{\partial f}{\partial x}\right)^{2}+\left(\frac{\partial f}{\partial y}\right)^{2}=4\) ，作变换 \(\left\{\begin{array}{l}x=u v \\ y=\frac{1}{2}\left(u^{2}-v^{2}\right)\end{array}\right.\) 使得 \(f(x, y)\) 变为 \(g(u, v)\) ．

\begin{itemize}
  \item[（1）] 若 \(a\left(\frac{\partial g}{\partial u}\right)^{2}+b\left(\frac{\partial g}{\partial v}\right)^{2}=u^{2}+v^{2}\) ，求 \(a, b\) 。
  \item[（2）] 对（1）式求出的 \(a, b\) ，利用极坐标变换 \(u=r \cos \theta, v=r \sin \theta\) 将 \(g\) 转化为 \(r, \theta\) 的表达式．
\end{itemize}

\subsection{题目9}
1．讨论 \(\int_{0}^{1} \frac{\sin ^{2} x}{x^{p}\left(x^{p}+\sin x\right)} \mathrm{d} x\) 在 \(p\) 取什么范围时收敛。

\subsection{题目10}
2．解答如下问题：\\
（1）设 \(u_{n}(x)\) 是 \([a, b]\) 上的单调递增函数，且 \(\sum_{n=1}^{\infty} u_{n}(a)\) 与 \(\sum_{n=1}^{\infty} u_{n}(b)\) 绝对收敛，问 \(\sum_{n=1}^{\infty} u_{n}(x)\) 在 \([a, b]\)上是否绝对收敛？是否一致收敛？\\
（2）设 \(u_{n}(x)\) 是 \([a, b]\) 上的单调递增函数，且 \(\sum_{n=1}^{\infty} u_{n}(a)\) 与 \(\sum_{n=1}^{\infty} u_{n}(b)\) 收敛，问 \(\sum_{n=1}^{\infty} u_{n}(x)\) 在 \([a, b]\) ↓ 是一致收敛？

\subsection{题目11}
3．设函数 \(f(x)\) 在 \([0,1]\) 上二阶可导，且 \(f(0)=f(1)=0, \min _{0 \leq x \leq 1} f(x)=-1\) ，证明： \(\max _{0 \leq x \leq 1} f^{\prime \prime}(x) \geq 8\) ．

\subsection{题目12}
4．设 \(f(x)\) 在 \([0,1]\) 上恒正连续，\(m, M\) 是 \(f(x)\) 在此区间上的最小值和最大值，证明：

\[
1 \leq \int_{0}^{1} \frac{1}{f(x)} \mathrm{d} x \int_{0}^{1} f(x) \mathrm{d} x \leq \frac{(m+M)^{2}}{4 M m} .
\]

\subsection{题目13}
5．设函数 \(f(x)\) 满足对任意的 \(x \in(a, b)\) ，存在 \(L_{x}>0\) 及 \(\delta_{x}>0\) ，满足

\[
\left|f\left(x_{1}\right)-f\left(x_{2}\right)\right|<L_{x}\left|x_{1}-x_{2}\right|, \forall x_{1}, x_{2} \in\left(x-\delta_{x}, x+\delta_{x}\right) \cap(a, b) .
\]

证明：存在 \(L>0\) ，对任意的 \(x, y \in(a, b)\) ，成立 \(|f(x)-f(y)|<L|x-y|\) ．

\subsection{题目14}
6．证明 \(\int_{1}^{+\infty} \frac{\cos (x t)}{1+x^{2}} \mathrm{~d} x\) 在 \((-\infty,+\infty)\) 上一致收敛，并证明 \(\lim _{t \rightarrow+\infty} \int_{1}^{+\infty} \frac{\cos (x t)}{1+x^{2}} \mathrm{~d} x=0\) ．
\end{document}
